\documentclass[11pt,a4paper]{article}
\usepackage[margin=1.08in]{geometry}
\usepackage{amsmath,amssymb,amsthm,mathtools,booktabs,array,microtype}
\usepackage[T1]{fontenc}
\usepackage{xcolor,hyperref,fancyhdr}
\definecolor{ink}{HTML}{193B50}
\hypersetup{colorlinks=true,linkcolor=ink,citecolor=ink,urlcolor=ink}
\pagestyle{fancy}\fancyhf{}\fancyhead[L]{\small\textsc{Integral Hodge lifting and divisible Sha}}\fancyhead[R]{\small J. V. de Mello}\fancyfoot[C]{\thepage}\setlength{\headheight}{14pt}
\newtheorem{theorem}{Theorem}[section]\newtheorem{lemma}[theorem]{Lemma}\newtheorem{corollary}[theorem]{Corollary}\newtheorem{proposition}[theorem]{Proposition}
\theoremstyle{remark}\newtheorem{remark}[theorem]{Remark}
\newcommand{\Sha}{\mathop{\mathrm{Sha}}}\newcommand{\Pic}{\operatorname{Pic}}\newcommand{\Br}{\operatorname{Br}}\newcommand{\NS}{\operatorname{NS}}\newcommand{\Gal}{\operatorname{Gal}}\newcommand{\Hom}{\operatorname{Hom}}\newcommand{\Z}{\mathbb Z}\newcommand{\Q}{\mathbb Q}\newcommand{\F}{\mathbb F}\newcommand{\G}{\mathbb G}\newcommand{\coker}{\operatorname{coker}}\newcommand{\loc}{\operatorname{loc}}\newcommand{\et}{\mathrm{\acute et}}\newcommand{\Div}{\operatorname{Div}}\newcommand{\ord}{\operatorname{ord}}\newcommand{\im}{\operatorname{im}}\newcommand{\diag}{\operatorname{diag}}\newcommand{\rank}{\operatorname{rank}}\newcommand{\Sel}{\operatorname{Sel}}
\title{\vspace{-1.2em}\textbf{Integral Hodge Lifting and Divisible\protect\\ Tate--Shafarevich Obstructions}\protect\\[6pt]\large Finite-level Poincar\'e reciprocity, arithmetic cover lattices, and geometric rigidity}
\author{Jo\~ao V\'ictor de Mello}
\date{November 26, 2026}
\begin{document}\maketitle\thispagestyle{plain}
\begin{abstract}
Let $E/K$ be a non-CM elliptic curve over a number field, $p\ge5$, and $X$ an everywhere locally soluble torsor under $E^h$ whose coordinates generate a finite $p$-primary subgroup $H$ of $\Sha(E/K)$. We give a finite-level criterion for lifting mixed Poincar\'e classes on $X\times E$ to arithmetic \'etale cohomology with lifts algebraic at every completion. A coefficient vector $a$ admits such a lift modulo $p^m$ if and only if its associated torsor class $a\xi$ belongs to $p^m\Sha(E/K)[p^\infty]$. The proof combines descent of multiplication covers, the universal Kummer class, cup products and a trace projection. The quotient of the locally algebraic lifting lattice by the algebraic lifting lattice is exactly $(H\cap p^m S)/p^mH$, where $S=\Sha(E/K)[p^\infty]$. At infinite precision the quotient becomes $H\cap\Div(S)$. We relate these formulas to the integral Brauer--polarization index, quantitative isogeny rigidity and finite Selmer descent, obtaining finite arithmetic upper bounds on the divisible corank and an exact stabilization criterion. The general vanishing of the divisible radical remains open; the comparison identifies rather than assumes the missing arithmetic input.
\end{abstract}
\noindent\textbf{Keywords:} Tate--Shafarevich group; Poincar\'e bundle; integral Hodge classes; Brauer group; Kummer covers; arithmetic duality.\par
\noindent\textbf{MSC (2020):} 11G05, 14F20, 14K05, 14G25, 11G40.
\section{Introduction and principal identities}
The finiteness of the Tate--Shafarevich group of an elliptic curve is one of the outstanding problems in arithmetic geometry. An obstruction to a naive Hodge-theoretic approach is that every torsor under a fixed abelian variety has the same geometric integral Hodge structure. The arithmetic lattices of descended cycles, however, need not agree. The distinction between divisibility in the Weil--Ch\^atelet group and divisibility within $\Sha$ is essential; see Creutz \cite{creutz} and Harari--Szamuely \cite{hs,hs2}.

Our principal comparison uses locally algebraic arithmetic lifts rather than only geometric Hodge classes. Set $S=\Sha(E/K)[p^\infty]$, $D=\Div(S)$, and let $H$ be generated by the coordinates of the torsor class $\xi$. For a positive integer $m$ put $R_m=\Z/p^m\Z$. We establish the following identities (all hypotheses and maps are specified below):
\begin{align}
\mathcal L_m/\mathcal A_m&\simeq(H\cap p^mS)/p^mH,\label{mainfinite}\\
\mathcal L_\infty/\mathcal A_\infty&\simeq H\cap D,\label{maininfinite}\\
\frac{\mathcal J_p(X\times E)}{\mathcal J_p(X)}&=\#(H\cap D).\label{mainj}
\end{align}
Here $\mathcal L_m$ and $\mathcal A_m$ are submodules of $R_m^h$, and $\mathcal J_p$ is an index of Tate-module lattices. The first equality gives a finite-level test for the obstruction in the second. The third compares it with locally trivial Brauer groups. No conjecture on $\Sha$ is used.

These results are complemented by the independent finite-descent estimate
\begin{equation}\label{mainbound}
 d_p\le \dim_{\F_p}\bigl(p^m S[p^{m+1}]\bigr),\qquad d_p=\rank_{\Z_p}T_p S.
\end{equation}
The right side is determined by finite Selmer groups and Mordell--Weil data. The estimate becomes an equality at sufficiently large precision. This does not give a universal a priori precision at which the right side vanishes.

\section{Torsors, Poincar\'e classes and geometric rigidity}
Fix an embedding $K\hookrightarrow\mathbb C$ and let $A=E^h$. The class $\xi=(\xi_1,\dots,\xi_h)$ belongs to $\Sha(A/K)=\Sha(E/K)^h$. All coordinates are assumed $p$-primary. Let $X$ be the corresponding torsor and put
\[
 \lambda:\Z^h\longrightarrow H,\qquad a\longmapsto a\xi=\sum_i a_i\xi_i,\qquad L=\ker\lambda.
\]
The map has index $[\Z^h:L]=\#H$. Translation acts trivially on geometric cohomology. Since $E$ is non-CM, the integral Lefschetz $(1,1)$ lattice on $X_\mathbb C$ is $\operatorname{Sym}_h(\Z)$ under the product principal polarization. In particular all degree-two integral Hodge classes are geometrically algebraic.
\begin{proposition}[Arithmetic polarization lattice]\label{pol}
The lattice of line bundles defined over $K$ is
\[
 \NS_K(X)=\{B\in\operatorname{Sym}_h(\Z): B\Z^h\subseteq L\}.
\]
For $H$ with invariant factors $q_1\mid\cdots\mid q_h$ (including factors $1$),
\[
 \min_{\mathcal L\ \mathrm{ample}}\chi(X,\mathcal L)=\#H,\qquad
 \NS(\bar X)/\NS_K(X)\simeq\bigoplus_{j=1}^h(\Z/q_j\Z)^j.
\]
\end{proposition}
\begin{proof}
The component of the Picard scheme associated to $B$ has torsor class $B\xi$, up to sign, by the theorem of the square and twisting by translations. A rational point of that component descends to a line bundle: its obstruction in $\Br(K)$ vanishes locally since $X(K_v)\ne\varnothing$, and hence globally. This proves the lattice statement. Riemann--Roch on an abelian variety gives $\chi(X,\mathcal L)=\det B$. For $B>0$, the index of $B\Z^h$ is divisible by $\#H$. If $L=U\diag(q_i)\Z^h$ with $U\in\operatorname{GL}_h(\Z)$, the positive definite matrix $U\diag(q_i)U^t$ realizes equality. The congruence $B\mapsto U^{-1}BU^{-t}$ reduces the lattice quotient to symmetric matrices whose $(i,j)$ entry is divisible by $\operatorname{lcm}(q_i,q_j)$, yielding the stated invariant factors.
\end{proof}
For $Y=X\times E$, the mixed part of a geometric divisor is a column $a\in\Z^h$. Its descent requires exactly $a\xi=0$: apply Proposition \ref{pol} to the torsor $(\xi,0)$ and a symmetric block matrix with off-diagonal column $a$. In particular the mixed algebraic lattice is $L$. Write $\mathcal A_m=(L+p^m\Z^h)/p^m\Z^h$ and $\mathcal A_\infty=L\otimes\Z_p$.

\section{Locally algebraic lifting: precise formulation}
Let $N=p^m$, $M=E[N]$, $Y=X\times E$. Use the Weil pairing $e_N:M\otimes M\to\mu_N$. Fix the sign of the normalized Poincar\'e line bundle $\mathcal P$ on $E\times E$ so that its mixed class agrees with the cup product of the universal Kummer classes. The sign is irrelevant to the submodules below. The geometric mixed degree-two summand of $Y$ is the direct sum of $h$ copies of $\operatorname{End}_{R_m}(M)$; a vector $a\bmod N$ represents the scalar endomorphism $a_i I$ in component $i$.
Define $\mathcal L_m\subseteq R_m^h$ to consist of vectors $a\bmod N$ for which there is a class
\[
 z\in H^2_{\et}(Y,\mu_N)
\]
whose geometric mixed component is the Poincar\'e class of $a$, and whose restriction at every completion $K_v$ belongs to the image of $\Pic(Y_{K_v})/N$ under the Kummer cycle map. Pure geometric components of $z$ are not prescribed. This is a subgroup, and $\mathcal A_m\subseteq\mathcal L_m$. The local algebraicity condition cannot be dropped: it is what distinguishes division in $\Sha$ from division in the larger Weil--Ch\^atelet group.

\begin{theorem}[Finite-level Poincar\'e lifting criterion]\label{lift}
For every $m\ge1$,
\begin{equation}\label{criterion}
 a\bmod N\in\mathcal L_m\quad\Longleftrightarrow\quad a\xi\in N S.
\end{equation}
Here $N S$ denotes multiplication by $N$ on $S$. Consequently
\begin{equation}\label{quot}
 \mathcal L_m/\mathcal A_m\simeq (H\cap N S)/NH.
\end{equation}
\end{theorem}
The right side is a purely arithmetic division group. We now prove the theorem in two steps: first an exact descent criterion for covers, then a cohomological equivalence between covers and mixed Hodge lifts.
\subsection{Descent of multiplication covers}
Let $u_E\in H^1_{\et}(E,M)$ be the class of $[N]:E\to E$. Choose a geometric origin on $X$ and identify $\bar X$ with $E^h$. For $a\in\Z^h$, denote by $a:\bar X\to E_{\bar K}$ the linear map with those coefficients. Its pullback of $u_E$ is independent of the geometric origin, since translations of $E_{\bar K}$ lift under $[N]$.
\begin{lemma}[Locally Kummer covers]\label{cover}
There is a class $u\in H^1_{\et}(X,M)$ restricting geometrically to $a^*u_E$, and restricting at each $K_v$ to the Kummer class of a morphism $X_{K_v}\to E$, if and only if $a\xi\in N S$.
\end{lemma}
\begin{proof}
Let $\xi_g\in A(\bar K)$ be a cocycle defining the affine Galois action $x\mapsto gx+\xi_g$ on $\bar X$. Over $\bar K$ the required $M$-torsor is
\[
 Z_a=\{(x,z):[N]z=a(x)\}\longrightarrow\bar X.
\]
This description also works if the cover is disconnected. An equivariant descent datum must act by $(x,z)\mapsto(gx+\xi_g,gz+\eta_g)$ with $N\eta_g=a\xi_g$. Indeed, choose any division of $a\xi_g$; two lifts differ by an $M$-equivariant automorphism over the connected base $\bar X$, necessarily translation by a constant element of $M$. The cocycle condition on descent is exactly the cocycle condition on $\eta_g\in E(\bar K)$. Changing the identification of the geometric cover changes $\eta$ by a coboundary. Changing the origin of $X$ also changes its class by a coboundary, using a division of the image of the translation. Thus a descended cover determines a torsor class $[\eta]\in H^1(K,E)$ with $N[\eta]=a\xi$.

At a completion, choose a point of $X(K_v)$ as origin. If $[\eta]_v=0$, choose $t\in E(\overline{K_v})$ with $\eta_g=gt-t$ (or its negative according to the coordinate convention). Translating the variable $z$ makes the equation a pullback of $[N]$ by a $K_v$-morphism to $E$, with linear part $a$. Conversely, any local Kummer cover arising from a morphism whose linear part is $a$ has zero root class. If the linear part is only congruent to $a$ modulo $N$, subtract $N$ times a morphism with the difference as linear part; this does not change its Kummer class. Hence local Kummerness is equivalent to $[\eta]\in\Sha(E/K)$. Since $a\xi$ is $p$-primary, any such root is $p$-primary: its prime-to-$p$ component is killed by $N$, which is invertible on that component. This proves necessity. Conversely a root in $S$ yields the required descent datum, after adjusting a representative so that $N\eta_g=a\xi_g$ as cocycles. The local argument gives all required local Kummer conditions.
\end{proof}
\subsection{The universal cup product and trace projection}
On $E\times E$, the two pullbacks $u_1,u_2$ of $u_E$ satisfy
\begin{equation}\label{cup}
 c_1(\mathcal P)=e_{N,*}(u_1\smile u_2)\quad\text{in }H^2_{\et}(E\times E,\mu_N).
\end{equation}
To verify this integrally, the geometric equality is the perfect Weil pairing and the K\"unneth formula (cf. \cite{sz1,sz2}). The difference lies in the first Hochschild--Serre filtration. Simultaneous multiplication by $2$ acts on both classes by $4$, using the normalized Poincar\'e bundle and functoriality of the Kummer class. On the filtration graded pieces of total degree two with positive base degree the eigenvalues are $2$ and $1$. Since $4-2$ and $4-1$ are units modulo $N$ for $p\ge5$, that difference is zero. This reasoning also works over every completion. In particular, for a morphism $f$ to $E$ the cup product of its Kummer class and $u_E$ is the cycle class of $(f,\mathrm{id})^*\mathcal P$, including for translated morphisms.

\begin{lemma}[Trace projection]\label{trace}
Let $\pi:Y=X\times E\to X$ and $\pi_E:Y\to E$ be the projections. There is a homomorphism
\[
 J:H^2_{\et}(Y,\mu_N)\longrightarrow H^1_{\et}(X,M),\qquad
 z\longmapsto\pm\pi_*(z\smile\pi_E^*u_E),
\]
where $\pi_*$ is the proper smooth trace and the sign is chosen compatibly with \eqref{cup}. Geometrically it sends the mixed Poincar\'e class of $a$ to $a^*u_E$ and annihilates the pure degree-two components. Locally it sends the Kummer cycle class of any line bundle on $Y_{K_v}$ to a Kummer class of a morphism $X_{K_v}\to E$.
\end{lemma}
\begin{proof}
The geometric claim is contraction against the universal degree-one class of $E$; the perfect pairing on $H^1(E)$ identifies it with the identity on the mixed summand. The normalization can be checked by a symplectic basis of $M$. At a completion, identify $X_{K_v}$ with $A$ by a local point. The Picard group of $A\times E$ decomposes as the pullbacks of line bundles on the two factors and the mixed normalized Poincar\'e line bundles attached to homomorphisms $A\to E$. This follows from the theorem of the cube, or the see-saw principle after restriction along the two origins. For a line bundle pulled back from $A$, the trace is zero by the projection formula and degree. For a bundle pulled back from $E$, express its divisor as a sum of closed points; the projection formula sends its cycle class cupped with $u_E$ to the sum of corestrictions of Kummer classes of the points. This is the Kummer class of a constant rational point, since corestriction commutes with the Kummer map. For the mixed Poincar\'e bundle attached to a homomorphism $f$, equation \eqref{cup} and the projection formula give $J(c_1((f,\mathrm{id})^*\mathcal P))=f^*u_E$ (the trace pairing of the two copies of $u_E$ is the identity under the fixed sign). The sum of these Kummer classes is the Kummer class of the sum of the corresponding morphisms. This proves the local claim.
\end{proof}
\begin{proof}[Proof of Theorem \ref{lift}]
If $a\xi\in N S$, take $u$ from Lemma \ref{cover}. The class $z=e_{N,*}(\pi^*u\smile\pi_E^*u_E)$ has the required geometric mixed component and is locally algebraic by \eqref{cup}. Conversely apply Lemma \ref{trace} to any locally algebraic lift $z$. The resulting class $J(z)$ satisfies Lemma \ref{cover}, proving necessity. Since $N\Z^h$ maps into $NH$, the condition is independent of the integral representative of $a\bmod N$. The map $a\mapsto a\xi\bmod NH$ is onto $(H\cap N S)/NH$ on $\mathcal L_m$, with kernel precisely $(L+N\Z^h)/N\Z^h$. This proves \eqref{quot}.
\end{proof}

\section{Exact stabilization and a finite arithmetic bound}
Finite Selmer theory implies that $S$ is a cofinitely generated $\Z_p$-module in the discrete sense. Noncanonically, $S=D\oplus F$ where $D\simeq(\Q_p/\Z_p)^d$ is its maximal divisible subgroup and $F\simeq S/D$ is finite. Put $b=\log_p\exp(F)$, with $b=0$ for $F=0$; also let $e=\log_p\exp(H)$.
\begin{corollary}[Finite-level stabilization]\label{stable}
For all $m\ge\max(1,b,e)$,
\[
 \mathcal L_m/\mathcal A_m\simeq H\cap D.
\]
If $\mathcal L_\infty$ is the lattice of $p$-adic coefficient vectors admitting compatible locally algebraic arithmetic lifts, then
\[
 0\to L\otimes\Z_p\to\mathcal L_\infty\xrightarrow{\lambda} H\cap D\to0.
\]
In particular $[\Z_p^h:\mathcal L_\infty]=\#\bar H\le\#F$, where $\bar H$ is the image of $H$ in $F$.
\end{corollary}
\begin{proof}
For $m\ge b$, $p^m S=D$; for $m\ge e$, $p^mH=0$. Apply Theorem \ref{lift}. A compatible lift necessarily belongs at each precision to the stated lattice, so its vector maps into $\bigcap_m p^mS=D$. Conversely, a class $a\xi\in D$ admits a compatible sequence of roots in $D$. Represent these roots by torsors $Y_j$ under $E$ and choose equivariant morphisms $Y_{j+1}\to Y_j$ for multiplication by $p$. Choose an equivariant morphism $X\to Y_0$ for the integral representative of $a$ modulo $p^e$, and pull back the composite covers $Y_j\to Y_0$. Their compatible Kummer classes and cup products yield compatible lifts. Changing the representative by $p^e$ times a vector adds a globally algebraic mixed class with $\Z_p$ coefficients, since $p^e\Z_p^h\subset L\otimes\Z_p$. Local points on all the torsors ensure the local Kummer property. Thus the image is exactly the claimed lattice; taking indices proves the formula.
\end{proof}
The index of the locally algebraic lifting lattice is uniformly bounded, whereas the index of the algebraic lattice is $\#H$ and can only grow through the divisible contribution. The important direction of this comparison is that boundedness of the former does not imply boundedness of the latter.

\begin{theorem}[Finite Selmer upper bound and geometric certificate]\label{certificate}
Put $S_j=S[p^j]$, and let $r=\rank E(K)$. For $m\ge0$ set
\[
 U_m=p^mS_{m+1}\subset S_1,\qquad u_m=\dim_{\F_p}U_m.
\]
Then
\begin{equation}\label{counts}
 d\le u_m=\log_p\#S_{m+1}-\log_p\#S_m
 =d+\#\{i:a_i>m\},
\end{equation}
where $F\simeq\bigoplus_i\Z/p^{a_i}\Z$ and $S_0=0$. In particular $u_m$ is nonincreasing, equals $d$ for $m\ge b$, and is independently computable from finite Selmer groups and the Mordell--Weil rank and torsion. If $u_m=0$ then $S=S_m$ and $p^mS=0$. Equivalently, for $m\ge1$, $S$ is finite if and only if there is a precision $m$ at which no nontrivial order-$p$ torsor admits a locally algebraic mixed lift of its primitive Poincar\'e class modulo $p^m$.
\end{theorem}
\begin{proof}
Multiplication by $p^m$ on $S_{m+1}$ has kernel $S_m$. Apply the structure theorem for $S$ to count the quotient; this proves \eqref{counts}. The exact Kummer sequence gives
\[\#S_j=\frac{\#\Sel_{p^j}(E/K)}{p^{jr}\#E(K)[p^j]}.\]
Here the denominator is the order of $E(K)/p^jE(K)$. This expresses the upper bound through finite descent data, not an assumption on the divisible rank. If $u_m=0$, the structure formula proves the asserted finiteness and exponent. For an order-$p$ class $\xi$, Theorem \ref{lift} with $h=1$ and $a=1$ says exactly that $\xi\in p^mS\cap S_1=U_m$. This proves the equivalence.
\end{proof}
\begin{remark}
Equality of two successive positive values of $u_m$ is not a stabilization certificate: finite elementary divisors may persist for many more levels. A zero value is a certificate. No universal bound on the required precision is known here.
\end{remark}
\subsection{A single-level Cassels--Tate certificate}
For $m\ge1$ consider the pairing matrix $B_{1,m}$ of the restriction of the Cassels--Tate pairing to $S_1\times S_m$. Its values lie in $\frac1p\Z/\Z\simeq\F_p$; the second argument factors through $S_m/p S_m$. Such pairings can be evaluated on finite-level classes without assuming finiteness of $S$; see \cite{ps,ms}.
\begin{theorem}[Cross-level duality and a one-level certificate]\label{dualcert}
The left radical of $B_{1,m}$ is exactly $U_m$. Hence
\[
 d\le \dim_{\F_p} S_1-\rank_{\F_p} B_{1,m}=u_m.
\]
If $B_{1,m}$ has full row rank, then $p^mS=0$ and $S$ is finite. This certificate uses $S_m$ and the pairing, rather than the next Selmer group $S_{m+1}$.
\end{theorem}
\begin{proof}
The Cassels--Tate pairing descends to a perfect alternating pairing on the finite group $F=S/D$. The image of $S_m$ in $F$ is $F[p^m]$, since the sequence with divisible kernel splits as groups. The annihilator of $F[p^m]$ under the perfect pairing is $p^mF$: containment follows by bilinearity and equality follows by comparing orders. Thus the radical on $S_1$ is $S_1\cap p^m S=p^mS_{m+1}$. Apply Theorem \ref{certificate} to obtain the claimed bound and certificate.
\end{proof}
At $m=1$ this is the usual residual pairing bound. At larger $m$ the second argument detects finite elementary divisors which are invisible to the pairing restricted to order-$p$ classes. In particular, the left radicals of these cross-level pairing matrices equal the subgroups of order-$p$ torsors admitting the locally algebraic mixed lifts of Theorem \ref{lift}. This provides a second, independent finite arithmetic test for the geometric obstruction.

\section{Integral Brauer--Hodge comparison}
Define $\mathcal B(X)=\ker[\Br_1(X)\to\prod_v\Br(X_{K_v})]$. For $A=E^h$ let $D_X$ be the finite image of the Picard connecting map $\NS(\bar X)\to H^1(K,A^\vee)$, identified with $B\mapsto B\xi$ under the product principal polarization. Then $D_X\subset S^h$. The low-degree Hochschild--Serre sequence and global Brauer reciprocity give
\begin{equation}\label{brauer}
 0\to\mathcal B(X)[p^\infty]\to S^h/D_X\xrightarrow{\beta_X}\Q/\Z,
\end{equation}
where $\beta_X$ is the sum of local invariants of a representative (equivalently the Cassels--Tate pairing against $\xi$, up to sign). In fact $H^1(K,\NS(\bar X))=0$ since the lattice is trivial and torsion-free. The image of locally trivial Picard classes is $\Sha(A^\vee)/D_X$; each local Brauer restriction is constant by a local point. Subtraction of a global constant makes all local restrictions zero precisely when the sum of invariants is zero. The constants inject into the local product by global reciprocity. This proves \eqref{brauer}; see \cite{hs2,ps} for the duality framework. A positive-degree global zero-cycle on $X$ annihilates the image of $\beta_X$, by evaluation and reciprocity.
\begin{theorem}[Integral Brauer--Hodge index]\label{brauerindex}
There is a canonical exact sequence
\[
 0\to (T_pS)^h\to T_p\mathcal B(X)\to D_X\cap D^h\to0.
\]
Thus $\mathcal J_p(X):=[T_p\mathcal B(X):(T_pS)^h]=\#(D_X\cap D^h)$, and
\[
 \rank_{\Z_p}T_p\mathcal B(X)=hd.
\]
Moreover, if $\bar D_X$ denotes the image in $F^h$ and $\bar\xi$ the image of $\xi$, then
\[
 \#\bigl(\mathcal B(X)[p^\infty]/\Div(\mathcal B(X)[p^\infty])\bigr)
 =\frac{(\#F)^h}{\ord(\bar\xi)\#\bar D_X}.
\]
\end{theorem}
\begin{proof}
The divisible subgroup of $S^h/D_X$ is the image of $D^h$. It lies in the kernel of the finite-image map $\beta_X$. Hence the divisible subgroup of $\mathcal B(X)[p^\infty]$ is $D^h/(D_X\cap D^h)$. Taking Tate modules of the quotient by the finite subgroup gives the displayed exact sequence: multiplication by $p^j$ and the snake lemma give the boundary into $(D_X\cap D^h)/p^j(D_X\cap D^h)$, and passage to the inverse limit is exact for these finite-level systems. The reduced formula follows from the perfect Cassels--Tate pairing on $F$ (\cite{ps}) and from the fact that the orthogonal complement of $\bar\xi$ has index its order. The subgroup $\bar D_X$ lies in that orthogonal complement by well-definedness of $\beta_X$ on the quotient. Taking orders proves the result.
\end{proof}
\begin{corollary}[Stable index ratio]\label{ratio}
For every $t\ge0$,
\[
 D_{X\times E^t}=D_X\oplus H^{\oplus t},\qquad
 \mathcal J_p(X\times E^t)=\mathcal J_p(X)\#(H\cap D)^t.
\]
In particular the ratio at $t=1$ is the order of the stabilized lifting quotient of Corollary \ref{stable}.
\end{corollary}
\begin{proof}
The block Picard connecting map sends $\bigl(\begin{smallmatrix}B&C\\C^t&D'\end{smallmatrix}\bigr)$ to $(B\xi,C^t\xi)$. The two parts vary independently, and $C^t\xi$ takes every value in $H^t$. Intersect with the maximal divisible subgroup and apply Theorem \ref{brauerindex}.
\end{proof}
\subsection{The geometric Brauer component}
At odd $p$, geometric Kummer theory and the K\"unneth formula identify $\Br(\bar X)[p^j]$ with $\mathfrak{sl}_2(\Z/p^j\Z)^{\binom h2}$, with conjugation action. This follows from the description of mixed classes as endomorphisms of $E[p^j]$ and the splitting of scalar and trace-zero matrices (\cite{sz1,sz2}). If $p^{c_p}M_2(\Z_p)\subset\Z_p[\rho_{E,p}(G_K)]$ (here $\rho_{E,p}$ denotes the full $p$-adic action), invariance implies that $p^{c_p}$ times a trace-zero matrix is scalar, hence zero. Therefore the invariant geometric Brauer group has order at most $p^{3c_p\binom h2}$, uniformly in $X$, and vanishes if the residual image contains $\operatorname{SL}_2(\F_p)$. The scalar algebra index $c_p$ is finite by Serre's open image theorem \cite{serre}. This removes the transcendental part, but not the divisible arithmetic group in Theorem \ref{brauerindex}.

\section{Quantitative geometric consequences}
For completeness we recall the link to isogenies and height energy, specifying the relevant uniformity. Put $\mu(X)=\min_{\mathcal L\ \mathrm{ample}}\chi(X,\mathcal L)$. Define $\mathfrak I(X)$ as the minimum degree of a dominant $K$-morphism from $X$ to any abelian variety of dimension $h$. By rigidity such a morphism is geometrically a translation followed by an isogeny $\varphi:A\to B$, satisfying $\varphi_*\xi=0$.
\begin{proposition}[Isogeny rigidity]\label{iso}
For $C=\prod_\ell\ell^{2hc_\ell}$, with $c_\ell$ the matrix algebra index above at $\ell$, one has
\[
 \mu(X)^2/C\le\mathfrak I(X)\le\mu(X)^2.
\]
For a $p$-primary $H$ one may take $C=p^{2hc_p}$; in particular if $c_p=0$, equality $\mathfrak I(X)=\mu(X)^2$ holds.
\end{proposition}
\begin{proof}
A matrix of rows forming a basis of $L$ annihilates $\xi$ and defines an endomorphism of $E^h$ of degree $(\#H)^2$, proving the upper bound. Conversely, take the $\ell$-primary kernel $U$ of an isogeny killing $\xi$; the quotient by $U$ kills the $\ell$-primary component since the remaining isogeny has prime-to-$\ell$ degree. Enlarge $U$ to the subgroup $U^{\rm sat}$ generated by the full matrix algebra on $E[\ell^j]^h$ for sufficiently large $j$. Then $\ell^{c_\ell}U^{\rm sat}\subset U$, so $[U^{\rm sat}:U]\le\ell^{2hc_\ell}$. By the matrix units, $U^{\rm sat}=W\oplus W$ in a basis of $E[\ell^j]$, for a subgroup $W\subset(\Z/\ell^j\Z)^h$. A matrix whose rows span the integral annihilator lattice of $W$ has this kernel; its determinant in absolute value is $\#W$. It kills the primary torsor class, so its determinant is divisible by the primary order of $H$. This gives $v_\ell(\deg\varphi)\ge2v_\ell(\#H)-2hc_\ell$. Multiply over primes. The constant is finite since the matrix algebra is full at all but finitely many primes by \cite{serre}. Proposition \ref{pol} completes the proof.
\end{proof}
For independent torsor coordinates of orders $q_i$, fix a non-torsion $P\in E(K)$ and the product polarization on the target $E^h$. The normalized canonical-height energy of a descent matrix $M$ is $\operatorname{tr}(M^tM)$. Each column $i$ is nonzero and divisible by $q_i$, hence the minimum is $\sum_i q_i^2$, attained by $\diag(q_i)$. Adding split factors preserves $\mu(X)$ by Proposition \ref{pol}; the integral mixed Poincar\'e descent cokernel for $X\times E^t$ is $H^t$. Moreover the degree-two integral cycle-class defect on $X$ equals $\NS(\bar X)/\NS_K(X)$ at $p$ after tensoring with $\Z_p$: Faltings's endomorphism theorem \cite{faltings} identifies the mixed invariants with scalar endomorphisms, generated primitively by the Poincar\'e classes. A nonzero defect rules out an isomorphism of integral $\Z_p$ Chow motives with $E^h$ by functoriality of cycle classes.

Combining Corollary \ref{stable}, Corollary \ref{ratio} and Proposition \ref{iso} gives the exact identity and bounds
\begin{align*}
 \mu(X)&=\#\bar H\cdot\#(\mathcal L_\infty/\mathcal A_\infty)
 =\#\bar H\frac{\mathcal J_p(X\times E)}{\mathcal J_p(X)},\\
 \mu(X)&\le\#F\cdot\#(\mathcal L_\infty/\mathcal A_\infty),\qquad
 \mathfrak I(X)\le(\#F)^2\#(\mathcal L_\infty/\mathcal A_\infty)^2.
\end{align*}
At a precision $m$ satisfying Theorem \ref{certificate} with $u_m=0$, all $p$-primary torsors in all fixed dimensions have $\mu(X)\le\#S_m$ and $\mathfrak I(X)\le(\#S_m)^2$, with $\#S_m$ computed from the finite Selmer group. Without that certificate we still have unconditional boundedness for subgroups $H$ disjoint from $D$. Bounded-degree finiteness follows from $\#H\le\sqrt{C\deg f}$ and finiteness of each finite-level Tate--Shafarevich group.

\section{The obstruction to unrestricted boundedness}
Suppose hypothetically that $D\simeq(\Q_p/\Z_p)^d$ with $d>0$. Take $h=d$, $H_n=D[p^n]$ and compatible generators defining $X_n$. The above theorems give the following exact growth laws; no example with $d>0$ is asserted.
\begin{center}\small
\begin{tabular}{@{}p{.60\linewidth}p{.30\linewidth}@{}}\toprule
Invariant & Value \\\midrule
$\mu(X_n)$ & $p^{nd}$\\
$\mathfrak I(X_n)$, if $c_p=0$ & $p^{2nd}$\\
Minimum fixed-target height energy (positive rank) & $dp^{2n}$\\
$\#(\NS(\bar X_n)/\NS_K(X_n))$ & $p^{nd(d+1)/2}$\\
$\mathcal L_\infty$ & $\Z_p^d$\\
$\mathcal A_\infty$ & $p^n\Z_p^d$\\
$\#(\mathcal L_\infty/\mathcal A_\infty)$ & $p^{nd}$\\
$\mathcal J_p(X_n\times E)/\mathcal J_p(X_n)$ & $p^{nd}$\\
\bottomrule\end{tabular}
\end{center}
In particular a uniform bound on the lifting quotient, even restricted to $h=1$, is equivalent to $d=0$. All these torsors have identical geometric Hodge structures, and the primitive mixed classes admit compatible locally algebraic arithmetic lifts. Thus an argument which only proves existence of such lifts cannot show that the corresponding classes are algebraic over $K$. This is an explicit obstruction to a proof of finiteness via geometric Hodge theory, and does not preclude an arithmetic Hodge theory argument with additional global input.

One possible target is a uniform arithmetic integral cycle comparison: a bound on the exponent of $\mathcal L_\infty/\mathcal A_\infty$ obtained independently of $S$. It would suffice to prove it for a single torsor coordinate, but its general validity is exactly the remaining finiteness question. A bound at one prime alone would not control all other primary components of $\Sha$.

For elliptic curves over $\Q$ of positive rank, under the hypotheses of Corollary C of Burungale--Castella--Grossi--Skinner \cite{bcgs} (in particular good ordinary, non-anomalous reduction and the relevant rational cyclotomic main conjecture), the first nonzero Kato derivative order is the fine Selmer corank. The ordinary $p^\infty$-Selmer group has corank $r+d$. Its localization at $p$ has corank one, because a non-torsion rational point has nonzero $p$-adic elliptic logarithm and the local Kummer group over $\Q_p$ has corank one. Hence the fine Selmer corank is $r-1+d$. Thus a finite test with $u_m=0$ proves the rank-sensitive order equality in that setting. This is a numerical implication, not a new higher motivic explicit reciprocity law.

\section{Conclusion and scope}
The finite-level criterion of Theorem \ref{lift} supplies a two-sided arithmetic realization of the obstruction previously seen only through polarization lattices and Tate-module indices. The constructive direction uses locally trivial division torsors; the converse uses a trace projection to recover a locally Kummer cover from an arithmetic lift. The stabilized quotient gives the exact intersection with the maximal divisible subgroup. Finite Selmer comparisons and Cassels--Tate evaluations give independent, finite upper bounds on the corank, while geometric rigidity gives explicit complexity bounds in the certified cases.

The argument does not prove that the independent upper bounds vanish for every elliptic curve, nor does it establish the unrestricted geometric boundedness conjecture or a higher Kato reciprocity law. The proposed contribution is the explicit finite-level Poincar\'e lifting comparison and its relation to the earlier geometric invariants. The foundational ingredients are classical; historical priority for this particular combination has not been established. Further independent verification of the proofs and comparison with the literature on descent, locally trivial coverings and integral cycle maps are necessary before asserting novelty.
\begin{thebibliography}{99}\small
\bibitem{creutz} B. Creutz, Locally trivial torsors that are not Weil--Ch\^atelet divisible, \emph{Bull. London Math. Soc.} 45 (2013), 935--942. \href{https://doi.org/10.1112/blms/bdt019}{doi:10.1112/blms/bdt019}.
\bibitem{hs} D. Harari and T. Szamuely, Galois sections for abelianized fundamental groups, \emph{Math. Ann.} 344 (2009), 779--800. \href{https://arxiv.org/abs/0808.2556}{arXiv:0808.2556}.
\bibitem{hs2} D. Harari and T. Szamuely, Local-global principles for 1-motives, \emph{Duke Math. J.} 143 (2008), 531--557. \href{https://arxiv.org/abs/math/0703845}{arXiv:math/0703845}.
\bibitem{sz1} A. N. Skorobogatov and Yu. G. Zarhin, The Brauer group of Kummer surfaces and torsion of elliptic curves, \emph{J. reine angew. Math.} 666 (2012), 115--140. \href{https://arxiv.org/abs/0911.2261}{arXiv:0911.2261}.
\bibitem{sz2} A. N. Skorobogatov and Yu. G. Zarhin, The Brauer group and the Brauer--Manin set of products of varieties, \emph{J. Eur. Math. Soc.} 16 (2014), 749--768. \href{https://doi.org/10.4171/JEMS/445}{doi:10.4171/JEMS/445}.
\bibitem{ps} B. Poonen and M. Stoll, The Cassels--Tate pairing on polarized abelian varieties, \emph{Ann. of Math.} 150 (1999), 1109--1149. \href{https://annals.math.princeton.edu/1999/150-3/p10}{Journal page}.
\bibitem{ms} A. Morgan and A. Smith, The Cassels--Tate pairing for finite Galois modules. \href{https://arxiv.org/abs/2103.08530}{arXiv:2103.08530}.
\bibitem{serre} J.-P. Serre, Propri\'et\'es galoisiennes des points d'ordre fini des courbes elliptiques, \emph{Invent. Math.} 15 (1972), 259--331.
\bibitem{faltings} G. Faltings, Endlichkeitss\"atze f\"ur abelsche Variet\"aten \"uber Zahlk\"orpern, \emph{Invent. Math.} 73 (1983), 349--366.
\bibitem{bcgs} A. Burungale, F. Castella, G. Grossi and C. Skinner, Non-vanishing of Kolyvagin systems and Iwasawa theory. \href{https://arxiv.org/abs/2312.09301}{arXiv:2312.09301}, Theorem C and Corollary C.
\end{thebibliography}
\end{document}
